% TEMPLATE FOR A NEW CASE STUDY
%
% 1. Copy this file to sections/case-studies/<short-name>.tex (e.g. ericksen.tex).
% 2. Replace every <...> below, and delete the parts you do not need.
% 3. In main.tex, add  \input{sections/case-studies/<short-name>}  under \part{Case studies}.
%    That line is what puts it in the PDF and on the website (as notes/pde/case-<short-name>.html).
%
% Labels: start every label in this file with the short name, so they never clash with
% another case study: \label{p:case-<short-name>}, \label{eq:<short-name>-system}, ...
% Links to Part I: \pref{energy-estimates}{energy estimates} prints "energy estimates (§5.1)" in the
% PDF and is a link on the website. The id is the file name in sections/foundations/;
% \pref{galerkin}{} uses the page's own title as the text.
%
% Boxes: \begin{step}{\routine{}Step 1. <title>} for a standard tool used in the standard way,
%        \begin{stepnew}{\newidea{}Step 2. <title>} for what is new in the paper.

\chapter{<Title of the case study>}\label{p:case-<short-name>}
\status{draft}  % stub, draft or done
\blurb{<One line shown under the link on the website's PDE page: the paper, and what it proves.>}

\begin{byline}Based on <Authors>, \emph{<Title of the paper>}, <Journal> <volume> (<year>) <pages>. \href{https://doi.org/<doi>}{doi:<doi>}\end{byline}

<One paragraph: what system this is, and what the paper proves about it.>
Steps are tagged \routine{} when a standard tool is used in the standard way, and \newidea{} when the argument does something that cannot be taken off the shelf.

% The concepts this case study uses (ids = file names in sections/foundations/). This makes the box
% here and the "Used in" lists on those concept pages.
\conceptsused{energy-estimates, galerkin, <concept-file-name>}

\section{The model}

<Where the system comes from, the unknowns, and the domain.>

\begin{equation}\label{eq:<short-name>-system}
  \left\{
  \begin{aligned}
    &<\text{first equation}>,\\
    &<\text{second equation}>,
  \end{aligned}
  \right.
\end{equation}
with boundary and initial conditions <...>.

\subsection*{Difficulties to keep in mind}

<What makes this system harder than the standard ones: a bad coupling term, a degenerate coefficient, a missing estimate, ...>

\section{The main results}

\begin{thmbox}{Theorem A (<e.g. local strong solutions>)}
<Statement.>
\end{thmbox}

\begin{thmbox}{Theorem B (<e.g. global weak solutions>)}
<Statement.>
\end{thmbox}

\section{The energy structure}

<The energy E and the dissipation D, and which multipliers produce them.>

\begin{step}{\routine{}The energy identity}
<Which equation is multiplied by what, and which terms cancel.>
\end{step}

\begin{stepnew}{\newidea{}<What is non-standard about the energy, if anything>}
<...>
\end{stepnew}

\section{Map of the proof of Theorem A}

\begin{step}{\routine{}Step 1. <title>}
<What is proved, and which tool from Part I does it.>
\end{step}

\begin{step}{\routine{}Step 2. <title>}
<...>
\end{step}

\section{Map of the proof of Theorem B}

\begin{stepnew}{\newidea{}Step 1. <The approximate problem>}
<...>
\end{stepnew}

\begin{step}{\routine{}Step 2. <Uniform bounds>}
<...>
\end{step}

\begin{step}{\routine{}Step 3. <Passing to the limit>}
<...>
\end{step}

\section{Exercises}

\begin{exercise}
  <Statement.>
  \begin{hint}
    <Hint.>
  \end{hint}
  \begin{solution}
    <Solution.>
  \end{solution}
\end{exercise}

\section{Open questions}

\begin{itemize}
\item <...>
\end{itemize}

\section*{References}

\begin{enumerate}[label={[\arabic*]}, leftmargin=*]
\item <Authors>, <Title>, \emph{<Journal>} <volume> (<year>) <pages>.
\end{enumerate}
